\section{Local obstructions and denominators on cubic curves}
\label{sec:local}

The previous sections give both positive descriptions and a global negative
result. Two further tools make the interface between geometry and finite
arithmetic explicit: reduction modulo a supplied integer, and the special
denominator structure of points on an integral Weierstrass equation.

\subsection{Finite congruence certificates}

For $m\ge2$, call a residue vector
$\bar x=(\bar x_0,\ldots,\bar x_n)\in\{0,\ldots,m-1\}^{n+1}$
\emph{unimodular modulo $m$} when
$\gcd(m,\bar x_0,\ldots,\bar x_n)=1$. For prime $m$, this means exactly
that not all coordinates are zero. For composite $m$, nonzeroness alone
is insufficient.

\begin{theorem}[A checkable certificate of no rational point]
\label{thm:modular}
Let $F_1,\ldots,F_s$ be homogeneous integer polynomials. If none of the
unimodular residue vectors modulo $m$ satisfies all equations
$F_i(\bar x)\equiv0\pmod m$, their projective zero locus has no rational
point. The premise is a primitive recursive test of at most $m^{n+1}$
vectors, and $\iPRA$ proves the implication.
\end{theorem}
\begin{proof}
A rational point has a primitive integral representative $x$. Reduce every
coordinate to $\bar x_i$ in $[0,m-1]$. Any common divisor of $m$ and all
$\bar x_i$ divides every $x_i$, so it is one; hence the residue vector is
unimodular. Polynomial evaluation respects congruence: each product does,
by induction on its number of factors, and a finite sum does as well.
Thus $F_i(x)=0$ implies $F_i(\bar x)\equiv0\pmod m$, contradicting the
tested premise. This is an intuitionistic proof of a negation from an
assumed point, not an appeal to a converse local--global principle.
\end{proof}

The empty conic $X^2+Y^2=3Z^2$ of Section~\ref{sec:conics} illustrates
why the modulus matters. Modulo three it has the nonzero vectors
$(0,0,1)$ and $(0,0,2)$. Modulo nine it has no unimodular solution: reduction
modulo three forces $3\mid X,Y$, and the original congruence modulo nine
then forces $3\mid Z$. The exact companion reproduces both residue counts.
Testing just the affine chart $Z=1$ would not be a valid general method
for excluding projective points, since a rational point may reduce to a
different chart.

The theorem is a one-way certificate. Failure to find an obstruction among
the tested moduli supplies no global existence conclusion. There is no
unbounded search for a successful modulus in any result here.

\subsection{Denominators on integral Weierstrass equations}

Consider
\[
      E_{A,B}:\quad y^2=x^3+Ax+B,\qquad A,B\in\Z.
\]
If $4A^3+27B^2\ne0$, its projective closure
$Y^2Z=X^3+AXZ^2+BZ^3$ is nonsingular. This can be checked directly.
At infinity the equation forces $X=0$, and the derivative with respect to
$Z$ at $[0:1:0]$ is $1$. At an affine singular point, the derivatives give
$y=0$ and $3x^2+A=0$, while the equation gives $x^3+Ax+B=0$.
Elimination gives $A=-3x^2$, $B=2x^3$ and therefore
$4A^3+27B^2=0$, contrary to the hypothesis. These arguments are polynomial
identities over characteristic zero. We need no group law for the results
of this section. Indeed their denominator assertions remain valid even
when the cubic is singular.

\begin{theorem}[Square and cube denominators]\label{thm:denominators}
Every rational affine point on $E_{A,B}$ has a unique expression
\[
        x=\frac a{d^2},\qquad y=\frac b{d^3},\qquad
        d>0,\quad\gcd(a,d)=\gcd(b,d)=1,
\]
with $a,b\in\Z$. These integers satisfy
\begin{equation}\label{eq:elliptic-integral}
             b^2=a^3+Aad^4+Bd^6.
\end{equation}
Conversely, integers with these conditions give a rational point. The
conversion in both directions is primitive recursive.
\end{theorem}
\begin{proof}
Write $x=a/q$ and $y=c/e$ in reduced form, with $q,e>0$. Let
\[
             T=a^3+Aaq^2+Bq^3.
\]
Then $y^2=T/q^3$. Since $T\equiv a^3\pmod q$ and $\gcd(a,q)=1$,
we have $\gcd(T,q)=1$, hence $\gcd(T,q^3)=1$. Also
$c^2/e^2$ is a reduced fraction. Uniqueness of reduced fractions therefore
gives
\[
                         e^2=q^3,\qquad c^2=T.
\]
For clarity, the uniqueness assertion follows from cross multiplication
and Bezout: $u/v=u'/v'$ with both denominators positive and both fractions
coprime implies $v\mid v'$ and $v'\mid v$, then $v=v'$ and $u=u'$.
It remains valid when a numerator is zero, since its reduced denominator
must then be one.

In the factorization of $q,e$, every prime exponent satisfies $2v=3u$.
Thus $u$ is even and $v$ is three times its half. The finite factorization
of Lemma~\ref{lem:arithmetic} constructs $d>0$ with $q=d^2$, $e=d^3$.
Equivalently the algorithm may take $d=\lfloor\sqrt q\rfloor$, since the
proof has established that $q$ is a square. Put $b=c$. Coprimality and
\eqref{eq:elliptic-integral} now follow from the reduced fractions.
The converse is obtained by dividing the integer equation by $d^6$.
For uniqueness, the reduced denominator of $x$ is $d^2$, so $d$ is its
unique positive square root and both numerators are determined.
\end{proof}

\begin{theorem}[Complete elliptic search at a given $x$-height]
\label{thm:elliptic-search}
Let $T\ge1$. All rational affine points on $E_{A,B}$ with
$\Ht(x)\le T$ are obtained by the following finite procedure:
\begin{enumerate}
\item Visit $1\le d\le\lfloor\sqrt T\rfloor$ and $-T\le a\le T$
with $\gcd(a,d)=1$.
\item Compute $W=a^3+Aad^4+Bd^6$. If $W\ge0$, compute its integer
square root $b_0$. Retain $(a/d^2,\pm b_0/d^3)$ exactly when $b_0^2=W$,
counting the two signs only once if $b_0=0$.
\end{enumerate}
The list is duplicate-free and has at most
$2(2T+1)\lfloor\sqrt T\rfloor$ points. Adjoining $[0:1:0]$ accounts
for the only point at infinity. This construction and its correctness are
provable in $\iPRA$.
\end{theorem}
\begin{proof}
Every retained pair satisfies the equation by direct substitution.
Coprimality of $a,d$ implies that $a/d^2$ is reduced. It also forces
$\gcd(b_0,d)=1$: a common prime would divide $a^3$ after reducing
\eqref{eq:elliptic-integral} modulo that prime, and hence would divide
$a$, a contradiction. Thus the displayed pairs are canonical, and the
uniqueness in the denominator theorem excludes duplicates beyond the
explicit zero-sign case.

Conversely the denominator theorem puts every point of the required height
in the form $a/d^2,b/d^3$ with $|a|\le T$ and $d^2\le T$. That pair
$(a,d)$ is visited and its value $W=b^2$ passes the square test. The sign
of $b$ is included. This proves completeness and the count bound.
For an explicit bound on the square-root search,
\[
      0\le W\le(1+|A|+|B|)T^3
\]
whenever $W\ge0$. If $C=1+|A|+|B|$, its square root is at most $CT^2$,
since $CT^3\le C^2T^4$. All loops and stored integers are therefore
explicitly bounded. At infinity, $Z=0$ forces $X=0$ and leaves the single
projective point $[0:1:0]$.
\end{proof}

This makes the familiar denominator geometry an actual bounded algorithm.
It does not provide a global height bound for all points on an arbitrary
cubic. For $E_1$, that global problem was already settled by the descent
proof: the affine list is exactly $(-1,0),(0,0),(1,0)$. For
$y^2=x^3-36x$, the point $(12,36)$ gives the rational right triangle
$(3,4,5)$ through the formulas in Section~\ref{sec:descent}. Thus the
same equations support both obstructed and solvable examples.
